\documentclass[11pt,a4paper]{article}
\usepackage[utf8]{inputenc}
\usepackage[T1]{fontenc}
\usepackage[french]{babel}
\usepackage{amsmath,amssymb}
\usepackage[margin=2cm]{geometry}
\usepackage{enumitem}
\usepackage{hyperref}

\newcommand{\vect}[1]{\overrightarrow{#1}}
\newcommand{\R}{\mathbb{R}}
\newenvironment{rem}{\par\medskip\noindent\textit{Remarque :}\quad}{\par\medskip}
\newenvironment{hyp}{\par\medskip\noindent\textit{Hypoth\`ese retenue :}\quad}{\par\medskip}

\title{Corrections d\'etaill\'ees\\Recueil TERMINISH -- Tome 1\\Classe de Premi\`ere D}
\author{Corrig\'e r\'edig\'e pour les \'el\`eves}
\date{}

\begin{document}
\maketitle

\noindent\textbf{Comment lire ce document.} Chaque exercice est corrig\'e \'etape par \'etape.
Lorsque l'\'enonc\'e du recueil est incomplet, ambigu ou contradictoire, la correction le signale
par une \emph{remarque} et indique l'hypoth\`ese retenue. Ces points sont \`a v\'erifier aupr\`es du professeur.

\tableofcontents
\newpage

\section{Chapitre 1 : Situations d'\'evaluation}

%% ===============================================================
\subsection{Exercice 1 : Probl\`eme de terrains}

\textbf{Consigne 1.} On pose $X=\cos x$. L'\'equation $(E)$ devient
$2X^2+(\sqrt3-\sqrt2)X-\dfrac{\sqrt6}{2}=0$.

\medskip\noindent\textbf{Discriminant :}
$\Delta=(\sqrt3-\sqrt2)^2+4\cdot 2\cdot\dfrac{\sqrt6}{2}=5-2\sqrt6+4\sqrt6=5+2\sqrt6=(\sqrt3+\sqrt2)^2$.

\noindent Donc
\[
X=\frac{-(\sqrt3-\sqrt2)\pm(\sqrt3+\sqrt2)}{4}
\quad\Longrightarrow\quad
X_1=\frac{\sqrt2}{2},\qquad X_2=-\frac{\sqrt3}{2}.
\]
\noindent\textbf{Solutions dans $]-\pi;\pi]$ :}
\begin{itemize}
  \item $\cos x=\dfrac{\sqrt2}{2}\iff x=\dfrac{\pi}{4}$ ou $x=-\dfrac{\pi}{4}$ ;
  \item $\cos x=-\dfrac{\sqrt3}{2}\iff x=\dfrac{5\pi}{6}$ ou $x=-\dfrac{5\pi}{6}$.
\end{itemize}
Les quatre sommets, sur le cercle de rayon $R=20$ m, sont dans l'ordre :
$-\frac{5\pi}{6}$, $-\frac{\pi}{4}$, $\frac{\pi}{4}$, $\frac{5\pi}{6}$.

\medskip\noindent\textbf{Aire.} Pour un polygone inscrit dans un cercle de centre $O$, l'aire est la somme
des aires des triangles $OS_iS_{i+1}$ : $\mathcal A=\frac12R^2\sum\sin\theta_i$, o\`u $\theta_i$ est l'angle
au centre entre deux sommets cons\'ecutifs. Ici les angles valent $\frac{7\pi}{12}$, $\frac{\pi}{2}$,
$\frac{7\pi}{12}$, $\frac{\pi}{3}$ (leur somme vaut bien $2\pi$). Comme
$\sin\frac{7\pi}{12}=\sin105^\circ=\dfrac{\sqrt6+\sqrt2}{4}$ :
\[
\mathcal A=\frac{400}{2}\left(2\cdot\frac{\sqrt6+\sqrt2}{4}+1+\frac{\sqrt3}{2}\right)
=200+100\sqrt3+100\sqrt6+100\sqrt2\approx 759{,}58\ \text{m}^2 .
\]
On retrouve la valeur annonc\'ee : $759{,}575\ \text{m}^2$.

\medskip\noindent\textbf{Consigne 2.}

\emph{Dimensions de $T_2$ :} soient $L$ la longueur et $\ell$ la largeur. On a $L\ell=360$ et
$(L+6)(\ell+6)=630$. Donc $L\ell+6(L+\ell)+36=630$, soit $L+\ell=39$.
$L$ et $\ell$ sont racines de $t^2-39t+360=0$ : $\Delta=81$, $t=\frac{39\pm 9}{2}\in\{24;15\}$.
Le terrain mesure $24$ m sur $15$ m, son p\'erim\`etre vaut $2(24+15)=78$ m.

\emph{Cl\^oture :} deux rang\'ees de fil de fer font le tour du terrain, sauf l'ouverture de 3 m.
Chaque rang\'ee mesure $78-3=75$ m, donc la longueur totale de fil vaut $2\times 75=150$ m.
\begin{rem}
Le texte dit ``2 rang\'ees du fil de fer''. Nous retenons deux fils parall\`eles de $75$ m chacun.
\end{rem}

\begin{itemize}
  \item Co\^ut du fil : $150\times 1550=232\,500$ F ;
  \item Main d'\oe uvre : $25\,650$ F ; co\^ut total de la cl\^oture : $258\,150$ F.
  \item Vente du terrain $T_1$ : $759{,}575\times 8000=6\,076\,600$ F.
\end{itemize}
\textbf{Reste :} $6\,076\,600-258\,150=\mathbf{5\,818\,450\ F}$.

%% ===============================================================
\subsection{Exercice 2 : \'Elevage de lapins}

\textbf{Consigne 1.} Notons $a$ le nombre de sacs d'aliment A et $b$ celui de sacs B.
Contraintes de temps de travail :
\[
\begin{cases}
a+3b\le 240 & \text{(groupe 1)}\\
2a+b\le 250 & \text{(groupe 2)}\\
a\ge 0,\ b\ge 0.
\end{cases}
\]
Sommets de la r\'egion admissible : $(0;0)$, $(125;0)$, $(0;80)$ et l'intersection des droites
$a+3b=240$ et $2a+b=250$ : de la seconde, $b=250-2a$ ; on remplace dans la premi\`ere :
$a+750-6a=240\Rightarrow a=102$, $b=46$.
Le nombre total de sacs $a+b$ est maximal en $(102;46)$ : $a+b=148$.
\begin{hyp}un sac permet d'alimenter un lapin ; le nombre maximal de lapins est donc $148$.
Le co\^ut des sacs n'intervient pas dans cette question.\end{hyp}

\textbf{Consigne 2.} Le polygone a pour sommets les images de la solution de
$(E)$ : $4\sin^2x+2(\sqrt3-1)\sin x-\sqrt3=0$.
Posons $Y=\sin x$ : $\Delta=4(\sqrt3-1)^2+16\sqrt3=16+8\sqrt3=(2+2\sqrt3)^2$, d'o\`u
$Y=\frac{-2(\sqrt3-1)\pm(2+2\sqrt3)}{8}$, soit $Y=\frac12$ ou $Y=-\frac{\sqrt3}{2}$.
Sur $[0;2\pi[$ : $\sin x=\frac12\iff x\in\{\frac{\pi}{6};\frac{5\pi}{6}\}$ et
$\sin x=-\frac{\sqrt3}{2}\iff x\in\{\frac{4\pi}{3};\frac{5\pi}{3}\}$.

Les sommets sont $A(\frac{\sqrt3}{2};\frac12)$, $B(-\frac{\sqrt3}{2};\frac12)$, $C(-\frac12;-\frac{\sqrt3}{2})$,
$D(\frac12;-\frac{\sqrt3}{2})$. Le quadrilat\`ere $ABCD$ est un trap\`eze : $(AB)\parallel(CD)$, avec
$AB=\sqrt3$, $CD=1$ et une hauteur $h=\frac12+\frac{\sqrt3}{2}$.
\[
\mathcal A=\frac{(\sqrt3+1)}{2}\cdot\frac{(1+\sqrt3)}{2}=\frac{(\sqrt3+1)^2}{4}=\frac{2+\sqrt3}{2}\approx 1{,}87\ \text{unit\'es d'aire}.
\]
\begin{rem} L'\'enonc\'e ne pr\'ecise pas l'unit\'e de longueur ; l'aire est donn\'ee en unit\'es d'aire.\end{rem}

%% ===============================================================
\subsection{Exercice 3 : C\'er\'emonie de dot}

Soit $G$ le barycentre de $\{(A;2),(B;3)\}$. Pour tout $M$ :
$2\vect{MA}+3\vect{MB}=5\vect{MG}$ et $\vect{MA}-\vect{MB}=\vect{BA}$.
La condition devient $5\,MG=AB=10$, soit $MG=2$ : \textbf{$(C)$ est le cercle de centre $G$ et de rayon $2$ m}
(\'echelle $1$ cm $\to 1$ m ; $AG=\frac35AB=6$ m et $BG=4$ m).

\medskip\noindent\textbf{Consigne 1.} Posons $Y=\sin x$ dans $(E)$ :
$-4\cos^2x+2(\sqrt3-\sqrt2)\sin x+4-\sqrt6=0$ devient, avec $\cos^2x=1-Y^2$,
$4Y^2+2(\sqrt3-\sqrt2)Y-\sqrt6=0$. Le discriminant vaut $4(\sqrt3+\sqrt2)^2$, donc
$Y=\frac{-(\sqrt3-\sqrt2)\pm(\sqrt3+\sqrt2)}{4}$, soit $Y=\frac{\sqrt2}{2}$ ou $Y=-\frac{\sqrt3}{2}$.
Les solutions dans $]-\pi;\pi]$ sont $\frac{\pi}{4}$, $\frac{3\pi}{4}$, $-\frac{\pi}{3}$, $-\frac{2\pi}{3}$.

Les sommets sont sur $(C)$, de centre $G$ et de rayon $2$ :
\begin{hyp}les angles sont mesur\'es depuis le centre $G$ du cercle $(C)$.\end{hyp}
L'aire du quadrilat\`ere inscrit vaut
\[
\mathcal A=\frac12\cdot 2^2\left(\sin\frac{\pi}{3}+\sin\frac{7\pi}{12}+\sin\frac{\pi}{2}+\sin\frac{7\pi}{12}\right)
=2+\sqrt2+\sqrt3+\sqrt6\approx 7{,}596\ \text{m}^2 .
\]
Co\^ut des travaux : $5000\times 7{,}59575=37\,978{,}75$ F, soit environ $37\,979$ F.

\medskip\noindent\textbf{Consigne 2.} Le nombre d'invit\'es n'est pas indiqu\'e dans l'\'enonc\'e.
\begin{rem} Il manque l'effectif de chaque groupe (A, B, C). Si l'on note $n_A,n_B,n_C$ ces effectifs, le nombre total
d'invit\'es est $N=n_A+n_B+n_C$ et le co\^ut du repas est $7200n_A+6200n_B+4300n_C$.
Merci de v\'erifier l'\'enonc\'e original pour obtenir la valeur num\'erique.\end{rem}

%% ===============================================================
\subsection{Exercice 4 : Am\'enagement sportif}

\textbf{Stade de hand-ball.} La relation $2\sin\theta-1=0$ donne $\sin\theta=\frac12$, soit
$\theta=\frac{\pi}{6}$ ou $\theta=\frac{5\pi}{6}$ modulo $2\pi$. Ces deux valeurs ne fournissent que deux points.
\begin{rem} Il faut quatre sommets pour un quadrilat\`ere. L'\'enonc\'e est incomplet (la relation utilis\'ee
donne deux points seulement). Avec la relation de l'exercice~10 ($2\cos2\theta-1=0$), on obtient quatre points,
et le rectangle correspondant a une aire de $\sqrt3$ unit\'es (voir exercice~10).
\end{rem}
Pour le rectangle $ABCD$ de centre $O$ de rayon $1$ (unit\'e $=10$ m) :
$\mathcal A=\sqrt3\times 1=\sqrt3$ unit\'es, soit $100\sqrt3\approx 173{,}2$ m$^2$ (car $1\ \text{u.a.}=10^2\ \text{m}^2$).

\textbf{Stade de volley-ball.} La rotation $r$ de centre $H(1;5)$ et d'angle $\frac{\pi}{2}$ envoie $E$ sur $G$.
On a $\vect{HE}=(-4;-3)$ et la rotation d'angle $\frac{\pi}{2}$ envoie $(x;y)$ sur $(-y;x)$ :
$\vect{HG}=(3;-4)$, donc $G(4;1)$.
\[
\vect{EF}=(3;-4),\qquad \vect{HG}=(3;-4),\qquad \vect{FG}=(4;3).
\]
$\vect{EF}=\vect{HG}$ : $EFGH$ est un parall\'elogramme ; $\vect{EF}\cdot\vect{FG}=12-12=0$ et $EF=FG=5$ :
c'est un \textbf{carr\'e} de c\^ot\'e $5$ unit\'es, soit $50$ m. Aire $=25$ u.a. $=2500$ m$^2$.

\medskip\noindent\textbf{Consigne 2.} Le gazon synthétique co\^ute $3500$ F (par m$^2$, hypoth\`ese) :
le terrain de volley-ball co\^ute $3500\times2500=8\,750\,000$ F.
\begin{rem} L'\'enonc\'e ne pr\'ecise pas si le gazon concerne le volley-ball ou le hand-ball, ni l'unit\'e du prix.
La valeur $n$ solution de $4+\sqrt{x-2}=x$ est : $\sqrt{x-2}=x-4$, donc $x-2=(x-4)^2$, soit $x^2-9x+18=0$, d'o\`u $x\in\{3;6\}$.
Comme $x-4\ge 0$, $n=6$. Cette valeur n'appara\^it pas clairement dans la suite de l'\'enonc\'e.\end{rem}

\medskip\noindent\textbf{Forage.} Le $k$-i\`eme m\`etre co\^ute $u_k=1000+50(k-1)$ F : suite arithm\'etique de raison $50$.
\begin{itemize}
  \item Co\^ut du $50^{\grave e}$ m\`etre : $u_{50}=1000+49\times50=\mathbf{3450\ F}$.
  \item Co\^ut de $d$ m\`etres : $\displaystyle S_d=\frac{d}{2}\bigl(2\cdot1000+(d-1)50\bigr)=25d^2+975d$.
  \item On cherche $d$ tel que $25d^2+975d=519\,750$, soit $d^2+39d-20\,790=0$.
  Discriminant : $\Delta=1521+83\,160=84\,681=291^2$, donc $d=\frac{-39+291}{2}=126$.
\end{itemize}
\textbf{La profondeur est de $126$ m.}

%% ===============================================================
\subsection{Exercice 5 : Coop\'erative agricole}

Soit $G$ le barycentre de $\{(A;2),(B;3)\}$. Comme $AB=10$ : $AG=6$, $BG=4$ et
$2MA^2+3MB^2=5MG^2+2GA^2+3GB^2=5MG^2+2\cdot36+3\cdot16=5MG^2+120$.
La condition $2MA^2+3MB^2=2120$ donne $5MG^2=2000$, soit $MG=20$.
La zone est le disque de centre $G$ et de rayon $20$ m. Son aire :
\[
\mathcal A=\pi\cdot20^2=400\times\frac{22}{7}=\frac{8800}{7}\approx 1257{,}14\ \text{m}^2.
\]

\medskip\noindent\textbf{Consigne 1.} Co\^ut d'am\'enagement : $\frac{8800}{7}\times 3000=\frac{26\,400\,000}{7}\approx 3\,771\,429$ F.
Co\^ut des terrains : $942\,000+2\,374\,000+1\,004\,000=4\,320\,000$ F.
\textbf{Budget total :} $4\,320\,000+3\,771\,429\approx \mathbf{8\,091\,429\ F}$.

\medskip\noindent\textbf{Consigne 2.} L'effectif de la coop\'erative n'est pas donn\'e dans l'\'enonc\'e.
Avec $15\,000$ F par participant, la somme de $1\,000\,000$ F suffit si et seulement si l'effectif $n$ v\'erifie
$15\,000\,n\le 1\,000\,000$, soit $n\le 66$ (car $1\,000\,000/15\,000\approx 66{,}7$).
\begin{rem} Il faut conna\^itre l'effectif pour conclure. Le calcul ci-dessus donne le seuil.\end{rem}

%% ===============================================================
\subsection{Exercice 6 : Jardin potager et probabilit\'es}

\textbf{Consigne 1.} Soit $L$ la longueur (le long du mur) et $\ell$ la largeur. Aire : $L\ell=392$.
Le grillage couvre une longueur et deux largeurs : $L+2\ell$. Donc $L=\dfrac{392}{\ell}$ et
\[
f(\ell)=\frac{392}{\ell}+2\ell,\qquad f'(\ell)=-\frac{392}{\ell^2}+2=\frac{2(\ell^2-196)}{\ell^2}.
\]
$f'$ s'annule en $\ell=14$ (n\'egative avant, positive apr\`es) : minimum.
Alors $L=28$ m, et le grillage mesure $28+2\times14=56$ m.
\textbf{Dimensions :} $28$ m (le long du mur) $\times\ 14$ m. Co\^ut minimal : $56\times1500=\mathbf{84\,000\ F}$.

\medskip\noindent\textbf{Consigne 2.} Le sac contient $10$ boules : $4$ noires (1 \`a 4), $4$ rouges (1 \`a 4), $2$ blanches (1 \`a 2).
Gagner : les trois boules sont de la m\^eme couleur, \emph{ou} toutes num\'erot\'ees $2$.
Il y a trois boules num\'erot\'ees $2$ (une noire, une rouge, une blanche).

\medskip\emph{Tirage simultan\'e} : $\binom{10}{3}=120$ tirages \'equiprobables.
\begin{itemize}
  \item M\^eme couleur : noires $\binom43=4$, rouges $\binom43=4$, blanches $\binom23=0$ : $8$ tirages.
  \item Toutes num\'erot\'ees $2$ : $\binom33=1$ tirage.
  \item Aucun tirage n'est \`a la fois de m\^eme couleur et tous num\'erot\'es $2$ (les trois couleurs seraient diff\'erentes).
\end{itemize}
$P_1=\dfrac{8+1}{120}=\dfrac{9}{120}=\dfrac{3}{40}$.

\medskip\emph{Tirage successif sans remise} : $10\times9\times8=720$ tirages ordonn\'es \'equiprobables.
\begin{itemize}
  \item M\^eme couleur : noires $4\times3\times2=24$, rouges $24$, blanches $0$ : $48$ tirages.
  \item Toutes num\'erot\'ees $2$ : $3!=6$ tirages ordonn\'es.
\end{itemize}
$P_2=\dfrac{48+6}{720}=\dfrac{54}{720}=\dfrac{3}{40}$.

\textbf{Conclusion :} les deux modes de tirage donnent la m\^eme probabilit\'e $\frac{3}{40}=0{,}075$. Il n'y a
\textbf{aucun tirage plus judicieux} ; le fermier a autant de chances dans les deux cas.

%% ===============================================================
\subsection{Exercice 7 : Route et zone de stockage}

\textbf{Consigne 1 (support de la route).} $f(x)=x^3-3x-1$ d\'efinie sur $\R$.
\begin{itemize}
  \item $f(-2)=-8+6-1=-3$ : le point $V(-2;-3)$ est bien sur la courbe.
  \item $f(2)=8-6-1=1$ : le point $F(2;1)$ est bien sur la courbe.
\end{itemize}
La route est l'arc de la courbe $\mathcal C_f$ compris entre $V$ et $F$, c'est-\`a-dire pour $x\in[-2;2]$.
\begin{rem} La route est un arc d'une courbe cubique : on la construit point par point (tableau de valeurs) sur $[-2;2]$.\end{rem}

\medskip\noindent\textbf{Consigne 2 (zone de stockage).} Plaçons les villes dans un rep\`ere orthonormé d'origine $A$ :
$A(0;0)$, $B(10;0)$ (car $AB=10$ km). Le point $C(x;y)$ v\'erifie $x^2+y^2=64$ et $(x-10)^2+y^2=36$ ; on obtient $x=6{,}4$, $y=4{,}8$.
Donc $C(6{,}4;4{,}8)$. (On constate $AC^2+BC^2=100=AB^2$ : le triangle est rectangle en $C$.)

Pour minimiser les co\^uts de trafic, on cherche le barycentre $S$ des villes pond\'er\'ees par leur tonnage
$650$ (A), $980$ (B), $540$ (C), total $2170$ :
\[
x_S=\frac{650\cdot0+980\cdot10+540\cdot6{,}4}{2170}=\frac{13\,256}{2170}\approx 6{,}11,\qquad
y_S=\frac{540\cdot4{,}8}{2170}=\frac{2592}{2170}\approx 1{,}19.
\]
\textbf{La zone de stockage est au point $S(6{,}11\ ;\ 1{,}19)$ (en km).}
\begin{hyp}le crit\`ere est la minimisation de $\sum m_i\,MM_i^2$ (moindres carr\'es pond\'er\'es), ce qui donne le barycentre.
Le tonnage le plus \'elev\'e (ville $B$) attire bien $S$ vers elle.\end{hyp}

%% ===============================================================
\subsection{Exercice 8 : PME et pr\^et bancaire}

Co\^ut $C(x)=x^2+202\,500$, prix de vente $1500$ F, capacit\'e $x\le1300$.
Recette $R(x)=1500x$ ; b\'en\'efice $B(x)=-x^2+1500x-202\,500$.

\textbf{Consigne 1.}
\begin{itemize}
  \item B\'en\'efice positif : $-x^2+1500x-202\,500>0\iff x^2-1500x+202\,500<0$.
  Discriminant $\Delta=1\,500^2-4\times202\,500=2\,250\,000-810\,000=1\,440\,000$, $\sqrt\Delta=1\,200$.
  Racines $x=\frac{1500\pm1200}{2}\in\{150;1350\}$. Le b\'en\'efice est positif pour $150<x<1350$.
  \textbf{Il faut produire au moins $151$ articles} pour gagner de l'argent.
  \item Maximum : $B(x)=-(x-750)^2+360\,000$, donc le b\'en\'efice est maximal pour $x=750$ (\`a l'int\'erieur de $[0;1300]$) :
  \textbf{$B_{\max}=360\,000$ F par mois.}
\end{itemize}

\textbf{Consigne 2.} Montant \`a rembourser : $10\,000\,000\times(1{,}125)^5$ avec capitalisation annuelle compos\'ee.
$(1{,}125)^5\approx1{,}802032$, donc la somme due est d'environ $18\,020\,320$ F.
Le titre foncier est estim\'e \emph{environ} \`a $18\,000\,000$ F : il est donc \textbf{tout juste insuffisant}
(il manque environ $20\,320$ F) ; il faudrait une valeur d'au moins $18\,020\,320$ F pour que la coop\'erative accepte le pr\^et.
\begin{rem} Le r\'esultat est tr\`es sensible \`a la valeur ``environ $18$ millions''. Si l'on suppose que la valeur est
exactement $18\,000\,000$ F, la garantie est insuffisante.\end{rem}

%% ===============================================================
\subsection{Exercice 9 : H\'eritage et minimisation}

\textbf{Consigne 1.} Soient $x$ la longueur et $y$ la largeur, avec $xy=400$. Le p\'erim\`etre vaut
$P=2(x+y)=2\left(x+\dfrac{400}{x}\right)=\dfrac{2x^2+800}{x}$ (formule demand\'ee).

\textbf{Consigne 2.} $P(x)=2x+\dfrac{800}{x}$ pour $x>0$. Alors
\[
P'(x)=2-\frac{800}{x^2}=\frac{2(x^2-400)}{x^2}.
\]
$P'(x)<0$ pour $0<x<20$ et $P'(x)>0$ pour $x>20$ : $P$ admet un minimum en $x=20$.
Alors $y=\dfrac{400}{20}=20$ : le terrain est un \textbf{carr\'e de c\^ot\'e $20$ m}, et son p\'erim\`etre minimal vaut $P_{\min}=80$ m.
\textbf{Dossou avait raison.}

%% ===============================================================
\subsection{Exercice 10 : Terrain de handball et d\'ecoration}

\textbf{Terrain de handball.} On a $2\cos2\theta-1=0\iff\cos2\theta=\frac12$, soit $2\theta=\pm\frac{\pi}{3}+2k\pi$,
donc $\theta\in\left\{\frac{\pi}{6};\frac{5\pi}{6};\frac{7\pi}{6};\frac{11\pi}{6}\right\}$ dans $[0;2\pi[$.
Les quatre sommets sont $(\pm\frac{\sqrt3}{2};\pm\frac12)$ : c'est un \textbf{rectangle} de c\^ot\'es $\sqrt3$ et $1$ (unit\'e $=10$ m).
\[
\mathcal A=\sqrt3\ \text{u.a.}=100\sqrt3\ \text{m}^2\approx173{,}21\ \text{m}^2.
\]
Co\^ut du gazon : $6000\times100\sqrt3=600\,000\sqrt3\approx\mathbf{1\,039\,230}$ F.

\medskip\noindent\textbf{Prix d'une planche.} Notons $p$ (planche), $t$ (t\^ole), $c$ (chevron). On a :
\[
\begin{cases}
20p+15t+5c=134\,500\\
10p+25t+5c=170\,500\\
10p+20t+5c=143\,500
\end{cases}
\]
La deuxi\`eme moins la troisi\`eme : $5t=27\,000$, donc $t=5\,400$.
La premi\`ere moins la troisi\`eme : $10p-5t=-9\,000$, donc $10p=18\,000$, soit $p=1\,800$.
Enfin $5c=143\,500-18\,000-108\,000=17\,500$, donc $c=3\,500$.
\textbf{Une planche de bois co\^ute $1\,800$ F.}

\medskip\noindent\textbf{Consigne 2 (motif).} $f(x)=\dfrac{x^2+ax+b}{2x-4}$ sur $\R\setminus\{2\}$.
\begin{itemize}
  \item $A(0;\frac12)\in\mathcal C$ : $f(0)=\dfrac{b}{-4}=\dfrac12$, donc $b=-2$.
  \item Tangente horizontale en $0$ : $f'(0)=0$. Or $f'(x)=\dfrac{(2x+a)(2x-4)-2(x^2+ax+b)}{(2x-4)^2}$,
  donc $f'(0)=\dfrac{-4a-2b}{16}=0\iff a=-\dfrac{b}{2}=1$.
\end{itemize}
\textbf{Ainsi $a=1$, $b=-2$ et} $f(x)=\dfrac{x^2+x-2}{2x-4}=\dfrac{(x+2)(x-1)}{2(x-2)}$.

\emph{Trac\'e :} $f(x)=\dfrac{x+3}{2}+\dfrac{2}{x-2}$ (division euclidienne : $x^2+x-2=(x-2)(x+3)+4$).
\begin{itemize}
  \item Asymptote verticale : $x=2$ ;
  \item Asymptote oblique : $y=\dfrac{x}{2}+\dfrac32$ ;
  \item Points remarquables : $f(0)=\frac12$, $f(1)=0$, $f(-2)=0$.
\end{itemize}
Pour $x>2$ la courbe est au-dessus de l'asymptote oblique, pour $x<2$ en dessous.

%% ===============================================================
\subsection{Exercice 11 : Route et lampadaires (Bac I 2024)}

\textbf{Consigne 1 (support).} $f(x)=\dfrac{x^2-5x+2}{x+2}$, d\'efinie sur $\R\setminus\{-2\}$.
Division : $x^2-5x+2=(x+2)(x-7)+16$, donc $f(x)=x-7+\dfrac{16}{x+2}$.
\begin{itemize}
  \item Asymptote verticale : $x=-2$ ; asymptote oblique : $y=x-7$.
  \item $f(0)=1$ : $A(0;1)$ est sur la courbe ; $f(6)=\dfrac{8}{8}=1$ : $B(6;1)$ est sur la courbe.
\end{itemize}
La route est l'arc de $\mathcal C_f$ pour $x\in[0;6]$.

\medskip\noindent\textbf{Consigne 2 (lampadaires).} $(x^2-1)(x^2-8x+15)=0\iff x\in\{-1;1\}$ ou $x\in\{3;5\}$.
Sur $[0;6]$ : $x\in\{1;3;5\}$ (le point $x=-1$ est hors de la route).
\begin{itemize}
  \item $f(1)=\dfrac{1-5+2}{3}=-\dfrac23$ : lampadaire en $\left(1;-\dfrac23\right)$ ;
  \item $f(3)=\dfrac{9-15+2}{5}=-\dfrac45$ : lampadaire en $\left(3;-\dfrac45\right)$ ;
  \item $f(5)=\dfrac{25-25+2}{7}=\dfrac27$ : lampadaire en $\left(5;\dfrac27\right)$.
\end{itemize}

%% ===============================================================
\subsection{Exercice 12 : Parcelle et pagnes}

\textbf{Consigne 1 (superficie).} Posons $C=\cos2x$. Comme $\cos4x=2\cos^22x-1$, l'\'equation devient
$2C^2-1-7C+4=0$, soit $2C^2-7C+3=0$. Racines : $C=\dfrac{7\pm5}{4}\in\{3;\tfrac12\}$.
$\cos2x\le1$ : on retient $\cos2x=\frac12$, donc $2x=\pm\frac{\pi}{3}+2k\pi$, soit $x=\pm\frac{\pi}{6}+k\pi$.
Dans $]-\pi;\pi]$ : $x\in\left\{\frac{\pi}{6};-\frac{\pi}{6};\frac{5\pi}{6};-\frac{5\pi}{6}\right\}$.
Les sommets forment un \textbf{rectangle} de c\^ot\'es $\sqrt3$ et $1$ (unit\'e $=100$ m) :
\[
\mathcal A=\sqrt3\ \text{u.a.}=10\,000\sqrt3\ \text{m}^2\approx 17\,300\ \text{m}^2\quad(\sqrt3=1{,}73).
\]

\textbf{Consigne 2 (pagnes).} Co\^ut en millions : $P(x)=\dfrac{1}{30}x^3-0{,}1x^2-0{,}3x+2$, $x\ge0$ (dizaines de pagnes).
Budget : $1\,000\,000$ F $=1$ million. Il faut $P(x)\le1$.
\[
P'(x)=\frac{x^2}{10}-\frac{x}{5}-\frac{3}{10}=\frac{x^2-2x-3}{10}=\frac{(x-3)(x+1)}{10}.
\]
$P$ d\'ecro\^it sur $[0;3]$, cro\^it apr\`es : son minimum est $P(3)=\dfrac{27}{30}-0{,}9-0{,}9+2=1{,}1$.
\textbf{Le co\^ut minimal est $1\,100\,000$ F $>1\,000\,000$ F.} M. Dodji \textbf{ne peut pas acheter} les pagnes avec sa somme.

%% ===============================================================
\subsection{Exercice 13 : TOGOTECH et transistors MOS}

\textbf{Consigne 1 (ordinateurs).} En millions de FCFA : co\^ut $C(x)=1120+0{,}00007x^2$, recette $R(x)=0{,}7x$.
\begin{itemize}
  \item Seuil de rentabilit\'e : $R=C\iff0{,}00007x^2-0{,}7x+1120=0$. En multipliant par $\frac{1}{0{,}00007}$ :
  $x^2-10\,000x+16\,000\,000=0$, $\Delta=36\,000\,000$, racines $x=\frac{10\,000\pm6\,000}{2}\in\{2\,000;8\,000\}$.
  Pas de perte si $2\,000\le x\le 8\,000$ : \textbf{au moins $2\,000$ ordinateurs} par mois.
  \item B\'en\'efice $B(x)=-0{,}00007x^2+0{,}7x-1120$ : maximal en $x=\dfrac{0{,}7}{2\times0{,}00007}=5\,000$.
  $B(5\,000)=3\,500-1\,120-1\,750=630$ : \textbf{b\'en\'efice maximal de $630$ millions de FCFA} pour \textbf{$5\,000$ ordinateurs} par mois.
\end{itemize}

\textbf{Consigne 2 (transistors).} Les temps de production (en secondes) forment une suite arithm\'etique :
$u_1=180$, raison $2$, donc $u_k=180+2(k-1)$.
Le temps total pour $n$ composantes est $S_n=\dfrac{n}{2}(2\cdot180+2(n-1))=n(n+179)$.
Dur\'ee journali\`ere : $3\,h\,59\,min=239\ \text{min}=14\,340$ s.
On cherche le plus grand $n$ tel que $n^2+179n\le14\,340$ : $\Delta=89\,401=299^2$, racine positive $n=\frac{-179+299}{2}=60$.
\textbf{Capacit\'e journali\`ere : $60$ composantes par appareil} (on a exactement $S_{60}=60\times239=14\,340$ s).

%% ===============================================================
\subsection{Exercice 14 : Entreprise de jouets et piscine}

\textbf{Consigne 1 (b\'en\'efice).} Chaque employ\'e produit $t$ jouets par jour ; trois employ\'es : $3t$ jouets.
Recette $500\times3t=1\,500t$. Co\^ut $3C(t)=30t^2-600t+6\,000$.
\[
B(t)=1\,500t-30t^2+600t-6\,000=-30t^2+2\,100t-6\,000.
\]
Maximum en $t=\dfrac{2100}{60}=35$. $B(35)=-36\,750+73\,500-6\,000=\mathbf{30\,750\ F}$ par jour,
pour $t=35$ jouets par employ\'e ($105$ jouets au total).
\begin{hyp}les trois employ\'es fabriquent la m\^eme quantit\'e $t$.\end{hyp}

\textbf{Consigne 2 (piscine).} L'\'equation devient, avec $Y=\cos x$ : $-4(1-Y^2)+2(\sqrt3-\sqrt2)Y+4-\sqrt6=0$, soit
$4Y^2+2(\sqrt3-\sqrt2)Y-\sqrt6=0$. Comme dans l'exercice~3, $Y\in\{\frac{\sqrt2}{2};-\frac{\sqrt3}{2}\}$, donc
$x\in\left\{\pm\frac{\pi}{4};\pm\frac{5\pi}{6}\right\}$ : quadrilat\`ere inscrit dans le cercle de rayon $R=5$ m.
\[
\mathcal A=\frac12R^2\left(2\sin105^\circ+1+\sin\frac{\pi}{3}\right)=\frac{25}{2}\left(1+\frac{\sqrt2+\sqrt3+\sqrt6}{2}\right)\approx 47{,}47\ \text{m}^2.
\]

\emph{Chiffre d'affaires de f\'evrier.} Posons $u=1+\frac{t}{100}$. Le CA de mars vaut $200\,000\cdot u\cdot(u+0{,}1)=312\,000$, soit
$u^2+0{,}1u-1{,}56=0$. $\Delta=6{,}25=2{,}5^2$, donc $u=\frac{-0{,}1+2{,}5}{2}=1{,}2$, soit $t=20$.
Le CA de f\'evrier vaut $200\,000\times1{,}2=\mathbf{240\,000\ F}$ (puis $240\,000\times1{,}3=312\,000$ en mars).

%% ===============================================================
\subsection{Exercice 15 : Taxi, m\'edicament et dispensaire}

\textbf{Consigne 1 (taxi).} Distance une fois : $10$ km.
\begin{itemize}
  \item Option 1 : suite g\'eom\'etrique $u_1=200$, $q=0{,}9$, $10$ termes :
  $\displaystyle S_1=200\,\frac{1-0{,}9^{10}}{1-0{,}9}\approx1\,302{,}64$ F.
  \item Option 2 : suite arithm\'etique $u_1=200$, raison $-20$, dernier terme $u_{10}=20$ :
  $\displaystyle S_2=\frac{10(200+20)}{2}=1\,100$ F.
\end{itemize}
Aller-retour : option 1 $\approx2\times1\,302{,}64=2\,605{,}29$ F ; option 2 $=2\,200$ F.
\textbf{Somme \'economis\'ee : environ $405{,}29$ F} avec la deuxi\`eme option.

\medskip\noindent\textbf{Heure du second m\'edicament.} Les mesures sont $(0;9{,}9)$, $(2;7{,}5)$, $(4;5{,}5)$, $(6;3{,}9)$, $(8;3)$.
Un ajustement exponentiel $y=a\,e^{kx}$ (r\'egression sur $\ln y$) donne $a\approx10$ et $k\approx-0{,}152$, soit $y\approx10\times0{,}86^x$.
\begin{rem} Le mod\`ele exponentiel ne s'annule jamais : il ne permet pas de d\'eterminer un instant d'\'elimination ``totale''.
L'\'enonc\'e ne donne pas de seuil d'\'elimination. Un ajustement polyn\^omial du second degr\'e
($y=0{,}05x^2-1{,}3x+9{,}9$) ne convient pas non plus (il ne s'annule pas et ne reproduit pas $y(8)=3$).
\textbf{La question est donc mal pos\'ee}. Il faut pr\'eciser le seuil ou le mod\`ele attendu.
\end{rem}

\medskip\noindent\textbf{Consigne 2 (dispensaire).} Populations : $A=1\,920$, $B=1\,280$, $C=640$, soit un total de $3\,840$.
L'\'enonc\'e ne fixe pas la position du dispensaire ; on utilise le crit\`ere du barycentre (voir ci-dessous).
\begin{hyp}le dispensaire est plac\'e au barycentre $D$ des trois villages pond\'er\'es par leur population :
$D=\mathrm{bar}\{(A;1920),(B;1280),(C;640)\}$, soit $\vect{AD}=\tfrac13\vect{AB}+\tfrac16\vect{AC}$.\end{hyp}
Les villages sont \'equidistants de $6$ km ($AB=AC=BC=6$, $\vect{AB}\cdot\vect{AC}=18$) :
\begin{itemize}
  \item $AD^2=\tfrac19\cdot36+\tfrac1{36}\cdot36+2\cdot\tfrac13\cdot\tfrac16\cdot18=4+1+2=7$, donc $AD=\sqrt7\approx2{,}65$ km ;
  \item $\vect{BD}=\vect{AD}-\vect{AB}=-\tfrac23\vect{AB}+\tfrac16\vect{AC}$, donc $BD^2=16+1-4=13$ et $BD=\sqrt{13}\approx3{,}61$ km ;
  \item $\vect{CD}=\vect{AD}-\vect{AC}=\tfrac13\vect{AB}-\tfrac56\vect{AC}$, donc $CD^2=4+25-20=9$ et $CD=3$ km.
\end{itemize}
Le village le plus peupl\'e $A$ est bien le plus proche du dispensaire ($\approx2{,}65$ km).

%% ===============================================================
\subsection{Exercice 16 : Entreprise de jouets et r\'emun\'erations}

\textbf{Consigne 1.} Le b\'en\'efice maximal est le m\^eme que celui de l'exercice~14 : \textbf{$30\,750$ F par jour}, pour $t=35$ jouets par employ\'e.

\medskip\noindent\textbf{Consigne 2.}
\emph{Chiffre d'affaires de d\'ecembre.} Octobre : $20\,000\,000$ F. Posons $u=1+\frac{t}{100}$ :
$20\,000\,000\cdot u(u+0{,}1)=23\,712\,000$, soit $u^2+0{,}1u-1{,}1856=0$.
$\Delta=0{,}01+4{,}7424=4{,}7524=2{,}18^2$, donc $u=\frac{-0{,}1+2{,}18}{2}=1{,}04$ : $t=4$.
Novembre : $20\,000\,000\times1{,}04=20\,800\,000$ F ; d\'ecembre : $20\,800\,000\times1{,}14=\mathbf{23\,712\,000\ F}$.

\emph{Manœuvre \`a r\'esilier.} On compte $10$ ann\'ees (de 2016 \`a 2025).
\begin{itemize}
  \item Abalo (arithm\'etique, premier terme $420\,000$, raison $20\,000$) :
  $\displaystyle S_A=\frac{10(2\cdot420\,000+9\cdot20\,000)}{2}=5\,100\,000$ F.
  \item Koudakpo (g\'eom\'etrique, $q=1{,}05$) :
  $\displaystyle S_K=420\,000\cdot\frac{1{,}05^{10}-1}{0{,}05}\approx5\,282\,715$ F.
\end{itemize}
Koudakpo a le plus co\^ut\'e : \textbf{c'est son contrat qui doit \^etre r\'esili\'e}.

%% ===============================================================
\subsection{Exercice 17 : Parc de loisir et stade de handball}

\textbf{Consigne 1 (d\'elai avant le pr\^et).} $A(\lambda)=3\lambda^2-24\lambda+96$.
$A$ est minimal en $\lambda=4$ ($A(4)=48$) : le chiffre d'affaires baisse puis remonte.
Pour $A(\lambda)=60$ : $3\lambda^2-24\lambda+36=0\iff\lambda^2-8\lambda+12=0$, d'o\`u $\lambda\in\{2;6\}$.
Le chiffre d'affaires d\'epasse $60$ millions pour la premi\`ere fois apr\`es la chute en $\lambda=6$ ans.
\textbf{Le pr\^et sera accord\'e au bout de $6$ ans}, soit $3$ ans apr\`es le constat de la chute.

\medskip\noindent\textbf{Consigne 2.}
\begin{rem} L'\'equation $(E)$ donn\'ee dans l'\'enonc\'e,
$4\cos^3x-2\sqrt3(\sqrt3-1)\sin2x-(\sqrt3+2)\sin x+\sqrt3=0$, ne se factorise pas sous la forme
$(\sin x+1)\bigl(4\sin^2x-2(\sqrt3+1)\sin x+\sqrt3\bigr)=0$ : en effet $x=\frac{3\pi}{2}$ annule le second membre
de la factorisation mais pas $(E)$ (on trouve $2+2\sqrt3\neq0$). L'\'enonc\'e doit \^etre v\'erifi\'e.
\end{rem}
En admettant la factorisation propos\'ee (sous la forme demand\'ee), les solutions sont :
\begin{itemize}
  \item $\sin x=-1\iff x=\frac{3\pi}{2}$ ;
  \item $4\sin^2x-2(\sqrt3+1)\sin x+\sqrt3=0$ : $\Delta=16-8\sqrt3=(2\sqrt3-2)^2$, donc $\sin x\in\{\frac{\sqrt3}{2};\frac12\}$,
  soit $x\in\{\frac{\pi}{3};\frac{2\pi}{3};\frac{\pi}{6};\frac{5\pi}{6}\}$.
\end{itemize}
Le budget (sable et ciment) ne peut pas \^etre calcul\'e sans conna\^itre le nombre de sommets et la surface du polygone,
qui d\'ependent de l'\'equation corrig\'ee. \textbf{Il faut v\'erifier l'\'enonc\'e.}

%% ===============================================================
\subsection{Exercice 18 : Transformations du plan}

\textbf{Donn\'ees.} $A(2;3)$, $B(-2;3)$, $C(-2;0)$, $G=\mathrm{bar}\{(A;2),(B;-3),(C;2)\}$ :
le poids total vaut $2-3+2=1$, donc $G=2A-3B+2C$ :
\[
x_G=2\cdot2-3\cdot(-2)+2\cdot(-2)=6,\qquad y_G=2\cdot3-3\cdot3+2\cdot0=-3.\qquad G(6;-3).
\]

\textbf{Consigne 1.} Pour tout $M$ : $2\vect{MA}-3\vect{MB}+2\vect{MC}=\vect{MG}$ (car $\sum w=1$).
L'\'egalit\'e $3\vect{MM'}=\vect{MG}$ donne $\vect{MM'}=\frac13\vect{MG}$, soit
$\vect{M'}=\frac23\vect{M}+\frac13\vect{G}$ en coordonn\'ees :
\[
\begin{cases}
x'=\frac23x+\frac13\cdot6=\frac23x+2\\[2pt]
y'=\frac23y+\frac13\cdot(-3)=\frac23y-1.
\end{cases}
\]
C'est bien l'expression annonc\'ee.

\textbf{Consigne 2.} Hypoth\`ese : $h$ est l'homoth\'etie de centre $G$ et de rapport $-3$, et $f$ l'application pr\'ec\'edente,
de centre $G$ et de rapport $\frac23$. Alors $f\circ h$ est l'homoth\'etie de centre $G$ de rapport $\frac23\times(-3)=-2$.
\begin{itemize}
  \item \textbf{Cercle $(C')$ :} centre $\Omega'=h'(B)=G-2\vect{GB}$. $\vect{GB}=(-8;6)$ donc $\Omega'=(6+16;\,-3-12)=(22;-15)$.
  Rayon $r'=|-2|\times2=4$. \textbf{Cercle de centre $(22;-15)$ et de rayon $4$.}
  \item \textbf{Droite $(D')$ :} $x'=-2x+18$ et $y'=-2y-9$. On exprime $x=\frac{18-x'}{2}$ et $y=\frac{-9-y'}{2}$,
  puis on remplace dans $y=-2x+3$ :
  $\dfrac{-9-y'}{2}=-(18-x')+3=x'-15$, donc $-9-y'=2x'-30$, soit $y'=-2x'+21$.
  \textbf{Droite $(D')$ : $2x+y-21=0$.}
\end{itemize}
\begin{rem} Le texte est partiellement illisible (``$h=f\circ h$''). L'interpr\'etation ci-dessus est la plus coh\'erente.
V\'erification : le point $(0;3)\in(D)$ a pour image $(18;-15)$, et $2\times18-15-21=0$.\end{rem}

%% ===============================================================
\subsection{Exercice 19 : Couloir et transformations}

\textbf{Consigne 1.} La figure n'est pas fournie ; la distance $AB$ ne peut donc pas \^etre exprim\'ee en fonction de $\alpha$.
\begin{rem} Sans la figure ni la relation liant $AB$ et $\alpha$, on ne peut ni trancher entre $GBEZE$ ($\frac{2\pi}{9}$ et $\frac{\pi}{3}$) et
$AMATSI$ ($\frac{\pi}{3}$ seul), ni valider leurs affirmations. Merci de fournir la figure.\end{rem}

\textbf{Consigne 2.}
\begin{itemize}
  \item \textbf{$T_1$ :} $x'=-\frac32x+\frac52$, $y'=-\frac32y+\frac52$. Point invariant : $x=-\frac32x+\frac52\Rightarrow x=1$ ;
  de m\^eme $y=1$. \textbf{$T_1$ est l'homoth\'etie de centre $\Omega(1;1)$ et de rapport $-\frac32$.}
  \item \textbf{$T_2$ :} matrice $\begin{pmatrix}-\frac{\sqrt3}{2}&-\frac12\\ \frac12&-\frac{\sqrt3}{2}\end{pmatrix}$.
  Colonnes orthogonales et normes $1$ : c'est une \textbf{rotation} de centre $O(0;0)$ (pas de translation).
  Angle $\theta$ : $\cos\theta=-\frac{\sqrt3}{2}$ et $\sin\theta=\frac12$, donc \textbf{$\theta=\frac{5\pi}{6}$}.
  \item \textbf{$T_3$ :} $G=\mathrm{bar}\{(A;1),(B;-2),(C;2)\}$, poids total $1$ : $G=A-2B+2C$, soit
  $x_G=1-2(-2)+2(-2)=1$ et $y_G=4-2\cdot1+2\cdot0=2$.
  \textbf{$G(1;2)$}. On a $-2\vect{MM'}=\vect{MG}$, donc $\vect{MM'}=-\frac12\vect{MG}$ et $\vect{GM'}=\frac32\vect{GM}$.
  \textbf{$T_3$ est l'homoth\'etie de centre $G(1;2)$ et de rapport $\frac32$.}
\end{itemize}
\textbf{Conclusion :} $T_1$ et $T_3$ ont des centres et des rapports diff\'erents ($(1;1)$, $-\frac32$ et $(1;2)$, $\frac32$) :
\textbf{LAWSON a tort}. Seul $T_2$ est une isom\'etrie, c'est une rotation d'angle $\frac{5\pi}{6}$.

\section{Chapitre 2 : Fonctions et limites}

%% ===============================================================
\subsection{Exercice 20 : \'Etude de fonction}

\textbf{1.(a)} Le d\'enominateur s'annule en $x=2$ : $\mathcal D_f=\R\setminus\{2\}$.

\textbf{1.(b)} Limites :
\begin{itemize}
  \item En $\pm\infty$ : $f(x)\sim\dfrac{2x^2}{x}=2x$, donc $\lim_{x\to+\infty}f=+\infty$ et $\lim_{x\to-\infty}f=-\infty$.
  \item En $2$ : le num\'erateur vaut $2\cdot4-14+7=1>0$. Donc $\lim_{x\to2^-}f=-\infty$ et $\lim_{x\to2^+}f=+\infty$.
\end{itemize}

\textbf{2.(a)} Division euclidienne : $2x^2-7x+7=(x-2)(2x-3)+1$. Donc
\[
f(x)=2x-3+\frac{1}{x-2} : \qquad a=2,\quad b=-3,\quad c=1.
\]
\textbf{2.(b)} $f(x)-(2x-3)=\dfrac{1}{x-2}\xrightarrow[x\to\pm\infty]{}0$ : la droite $(D):y=2x-3$ est asymptote oblique.

\textbf{2.(c)} $\lim_{x\to2}f(x)=\pm\infty$ : la droite $x=2$ est l'autre asymptote (verticale).

\textbf{3.} Soit $h\neq0$. $f(2+h)=2(2+h)-3+\dfrac1h=1+2h+\dfrac1h$, et
$f(2-h)=1-2h-\dfrac1h$. Donc $f(2+h)+f(2-h)=2$ : le point $A(2;1)$ est un \textbf{centre de sym\'etrie} de $(C)$.

%% ===============================================================
\subsection{Exercice 21 : Fonction rationnelle}

\textbf{1.} $\mathcal D_f=\R\setminus\{1\}$ ; les limites aux bornes se lisent sur le tableau de variations
(non reproduit dans le texte extrait) : elles sont $\pm\infty$ de part et d'autre de $x=1$, et la droite
$(D)$ donne le comportement en $\pm\infty$.

\textbf{2.} Avec $f(x)=\dfrac{ax^2+bx+c}{x-1}$, les coefficients se d\'eduisent du tableau de variations
(``on suppose'' : $a,b,c$ sont fix\'es par les donn\'ees du tableau).
\begin{rem} Le tableau n'est pas reproduit dans le texte. Dans la suite, on prend $a=1$, $b=-4$, $c=7$,
ce qui est coh\'erent avec la question~3.\end{rem}

\textbf{3.} On pose $f(x)=\dfrac{x^2-4x+7}{x-1}$.
\begin{enumerate}[label=(\alph*)]
  \item Division : $x^2-4x+7=(x-1)(x-3)+4$, donc $f(x)=x-3+\dfrac{4}{x-1}$.
  Comme $\dfrac{4}{x-1}\to0$ en $\pm\infty$, la droite $(D):y=x-3$ est asymptote \`a $(C_f)$.
  \item Position relative : $f(x)-(x-3)=\dfrac{4}{x-1}$.
  \begin{itemize}
    \item Si $x>1$, $f(x)-(x-3)>0$ : $(C_f)$ est \textbf{au-dessus} de $(D)$.
    \item Si $x<1$, $f(x)-(x-3)<0$ : $(C_f)$ est \textbf{en dessous} de $(D)$.
  \end{itemize}
  \item Axes : en $x=0$, $f(0)=\dfrac{7}{-1}=-7$ : point $(0;-7)$. Avec l'axe des abscisses, $x^2-4x+7=0$ a un discriminant
  $16-28<0$ : \textbf{pas de point d'intersection} avec $(Ox)$.
  \item Centre de sym\'etrie : $f(1+h)+f(1-h)=(1+h-3+\tfrac4h)+(1-h-3-\tfrac4h)=-4$.
  Donc $I(1;-2)$ est centre de sym\'etrie, puisque $\dfrac{-4}{2}=-2$ est la moyenne des ordonn\'ees.
\end{enumerate}

\textbf{4.} \emph{Trac\'e :} on place $(D):y=x-3$ (passant par $(3;0)$ et $(0;-3)$), l'asymptote verticale $x=1$, le point $(0;-7)$,
et le centre $I(1;-2)$. La branche gauche passe sous $(D)$, la branche droite au-dessus.

%% ===============================================================
\subsection{Exercice 22 : Centre de sym\'etrie}

\textbf{1.(a)} Pour $x\neq4$ : $h(x)=\dfrac{3x+5}{x-4}=3+\dfrac{17}{x-4}$. Donc $h(4+t)+h(4-t)=3+\frac{17}{t}+3-\frac{17}{t}=6$ :
$\Omega(4;3)$ est centre de sym\'etrie de $(C_h)$.

\textbf{1.(b)} Pour $y\neq3$ : $y=\dfrac{3x+5}{x-4}\iff y(x-4)=3x+5\iff x(y-3)=4y+5\iff x=\dfrac{4y+5}{y-3}$.
Ainsi $h$ est bijective de $\R\setminus\{4\}$ sur $\R\setminus\{3\}$ et
$\boxed{h^{-1}(x)=\dfrac{4x+5}{x-3}}$.

\textbf{2.(a)} $g(x)=3x^2-6x$ : $g(-x)=3x^2+6x$, qui n'est ni $g(x)$ ni $-g(x)$. \textbf{$g$ n'est ni paire ni impaire.}

\textbf{2.(b)} $g(1+t)=3(1+2t+t^2)-6(1+t)=3t^2-3$, qui est une fonction paire de $t$. Donc la droite $(\Delta):x=1$
est un \textbf{axe de sym\'etrie} de $(C_g)$.

\textbf{2.(c)} $\mathcal D_{f\circ g}=\{x\in\R : g(x)\neq4\}$ (en prenant $f=h$, l'\'enonc\'e \'etant peu lisible).
$3x^2-6x-4=0$ a pour racines $1\pm\dfrac{\sqrt{21}}{3}$ :
$\mathcal D_{f\circ g}=\R\setminus\left\{1-\tfrac{\sqrt{21}}{3};\,1+\tfrac{\sqrt{21}}{3}\right\}$.
Puis $(f\circ g)(x)=\dfrac{3g(x)+5}{g(x)-4}=\dfrac{9x^2-18x+5}{3x^2-6x-4}$.

\textbf{3.(a)} $\dfrac{x^2-3x+2}{x-1}=\dfrac{(x-1)(x-2)}{x-1}=x-2$ pour $x\neq1$ : la limite vaut $\mathbf{-1}$.

\textbf{3.(b)} $\dfrac{x^2-3x+2}{2-x}\sim\dfrac{x^2}{-x}=-x$ : $\lim_{x\to+\infty}=\mathbf{-\infty}$.

%% ===============================================================
\section{Chapitre 3 : Suites num\'eriques}

%% ===============================================================
\subsection{Exercice 23 : Suite d\'efinie par r\'ecurrence}

$U_0=0$, $U_{n+1}=\frac13U_n+3$.

\textbf{1.} $U_1=3$, $U_2=\dfrac{3}{3}+3=4$. (Les points $(0;0)$, $(1;3)$, $(2;4)$ se placent dans le rep\`ere $(O;I;J)$.)

\textbf{2.} Posons $V_n=U_n-\dfrac92$. Alors
\[
V_{n+1}=\tfrac13U_n+3-\tfrac92=\tfrac13U_n-\tfrac32=\tfrac13\left(U_n-\tfrac92\right)=\tfrac13V_n.
\]
$(V_n)$ est g\'eom\'etrique de raison $\frac13$ et de premier terme $V_0=-\frac92$.

\textbf{3.} $V_n=-\dfrac92\left(\dfrac13\right)^n$ et $U_n=\dfrac92-\dfrac92\left(\dfrac13\right)^n=\dfrac92\left(1-\left(\tfrac13\right)^n\right)$.

\textbf{4.(a)} $S_n=\sum_{k=0}^nV_k=-\dfrac92\cdot\dfrac{1-(1/3)^{n+1}}{1-1/3}=-\dfrac{27}{4}\left(1-\left(\tfrac13\right)^{n+1}\right)$.

\textbf{4.(b)} $S'_n=\sum_{k=0}^nU_k=\dfrac92(n+1)-\dfrac92\sum_{k=0}^n\left(\tfrac13\right)^k$
$=\dfrac92(n+1)-\dfrac{27}{4}\left(1-\left(\tfrac13\right)^{n+1}\right)$.
En d\'eveloppant : $S'_n=\dfrac92n+\dfrac92-\dfrac{27}{4}+\dfrac{9}{4}\left(\tfrac13\right)^n=\dfrac94\left[\left(\tfrac13\right)^n+2n-1\right]$.
\begin{rem} L'\'enonc\'e \'ecrit ``$-2n-1$''. Le bon signe est $+2n-1$ : pour $n=1$, $S'_1=U_0+U_1=3$ et
$\frac94\left(\frac13+2-1\right)=3$, alors que la formule avec $-2n$ donnerait $-6$.\end{rem}

%% ===============================================================
\subsection{Exercice 24 : Suites}

\textbf{1.(a)} $V_0=\dfrac{1}{U_0-3}=\dfrac{1}{4-3}=\mathbf{1}$.

\textbf{1.(b)} $U_{n+1}-3=\dfrac{4U_n-9-3U_n+6}{U_n-2}=\dfrac{U_n-3}{U_n-2}$. Donc
\[
V_{n+1}=\frac{U_n-2}{U_n-3}=\frac{(U_n-3)+1}{U_n-3}=1+V_n.
\]
$(V_n)$ est \textbf{arithm\'etique} de raison $1$ (et de premier terme $1$) : $V_n=n+1$.
Ainsi $U_n=3+\dfrac{1}{n+1}=\dfrac{3n+4}{n+1}$. (V\'erification : $U_1=\frac72=\frac{4\cdot4-9}{4-2}$.)

\textbf{2.} $(w_n)$ est g\'eom\'etrique de raison $q=\frac23$ et $w_2=5$.

\textbf{2.(a)} $w_n=w_2\,q^{\,n-2}=\mathbf{5\left(\tfrac23\right)^{n-2}}$ : r\'eponse \textbf{b)}.

\textbf{2.(b)} $S_n=w_2+\cdots+w_n$ compte $n-1$ termes :
$S_n=5\,\dfrac{1-(2/3)^{n-1}}{1-2/3}=15\left[1-\left(\tfrac23\right)^{n-1}\right]$ : r\'eponse \textbf{d)}.

\textbf{2.(c)} Comme $0<\frac23<1$, $\lim_{n\to+\infty}w_n=\mathbf{0}$ : r\'eponse \textbf{a)}.

%% ===============================================================
\section{Chapitre 4 : G\'eom\'etrie plane}

%% ===============================================================
\subsection{Exercice 25 : Transformations du plan}

$f:\ x'=-2x+3,\ y'=-2y+3$.

\textbf{1.} Point invariant : $x=-2x+3\iff x=1$ et $y=-2y+3\iff y=1$. \textbf{$\Omega(1;1)$} est l'unique point invariant.

\textbf{2.} Comme $\Omega$ est invariant : $\vect{\Omega M'}=-2\,\vect{\Omega M}$.

\textbf{3.} $f$ est l'\textbf{homoth\'etie de centre $\Omega(1;1)$ et de rapport $-2$}.

\textbf{4.} Pour inverser : $x'=-2x+3\Rightarrow x=\dfrac{3-x'}{2}$, et de m\^eme $y=\dfrac{3-y'}{2}$.
\[
f^{-1}:\quad x'=-\tfrac12x+\tfrac32,\qquad y'=-\tfrac12y+\tfrac32.
\]
C'est l'homoth\'etie de centre $\Omega(1;1)$ et de rapport $-\frac12$.

%% ===============================================================
\subsection{Exercice 26 : Homoth\'etie et barycentre}

\textbf{Partie A.} $h:\ \vect{AM'}=\tfrac32\vect{AM}+\tfrac12\vect{BA}$.

\textbf{1.(a)} $\vect{BM'}=\vect{BA}+\vect{AM'}=\vect{BA}+\tfrac32\vect{AM}+\tfrac12\vect{BA}
=\tfrac32\left(\vect{BA}+\vect{AM}\right)=\tfrac32\vect{BM}$.

\textbf{1.(b)} Ainsi $\vect{BM'}=\tfrac32\vect{BM}$ : $h$ est l'\textbf{homoth\'etie de centre $B$ et de rapport $\frac32$}.

\textbf{Partie B.} $ABC$ rectangle en $A$, $AB=4$ cm, $AC=6$ cm.

\textbf{2.(a)} Soit $G=\mathrm{bar}\{(A;1),(B;2)\}$ : $\vect{MA}+2\vect{MB}=3\vect{MG}$. La condition devient
$3\,MG\times MC=10$, soit $MG\times MC=\dfrac{10}{3}$.
\begin{rem} Cet ensemble n'est pas un cercle ni une droite (produit de deux distances). Il s'agit d'une courbe
de type ``Cassini'' (ensemble de points dont le produit des distances \`a deux points est constant).
L'\'enonc\'e attendait sans doute une autre forme ; il faut le v\'erifier.\end{rem}

\textbf{2.(b)} Par l'homoth\'etie $h$ (centre $B$, rapport $\frac32$), l'ensemble $(C')=h(C)$ est l'image de cette courbe.
\begin{rem} Conform\'ement \`a la remarque ci-dessus, $(C')$ est aussi une courbe de m\^eme type.\end{rem}

%% ===============================================================
\subsection{Exercice 27 : Barycentre et ensembles de points}

$ABC$ \'equilat\'eral de c\^ot\'e $a$ et de centre $I$. Soit $K$ le milieu de $[AB]$.

\textbf{1.} Pour tout $M$ : $\vect{MA}+\vect{MB}=2\vect{MK}$. La condition devient
\[
\vect{MC}\cdot2\vect{MK}=-MC^2\iff\vect{MC}\cdot\left(\vect{MC}+2\vect{MK}\right)=0.
\]
Dans le rep\`ere d'origine $K$ : on obtient $|M-\tfrac23C|^2=\tfrac19|C|^2$.
\textbf{$(E_1)$ est le cercle de centre $\Omega=\mathrm{bar}\{(A;1),(B;1),(C;4)\}$ et de rayon $\dfrac{a\sqrt3}{6}$.}
Il passe par $C$ et par $I$ (car $KC=a\frac{\sqrt3}{2}$ et $\Omega C=\frac13KC=\frac{a\sqrt3}{6}$).

\textbf{2.(a)} $f(M)=2MA^2+MB^2+MC^2$ :
\[
f(A)=0+a^2+a^2=2a^2,\qquad f(B)=2a^2+0+a^2=3a^2,\qquad f(C)=2a^2+a^2+0=3a^2.
\]
\textbf{2.(b)} $G=\mathrm{bar}\{(A;2),(B;1),(C;1)\}$ : $\vect{AG}=\tfrac14(\vect{AB}+\vect{AC})$.

\textbf{2.(c)} Dans le rep\`ere $A(0;0)$, $B(a;0)$, $C\left(\tfrac a2;\tfrac{a\sqrt3}{2}\right)$ :
$G\left(\tfrac{3a}{8};\tfrac{a\sqrt3}{8}\right)$. Alors
\[
AG^2=\frac{3a^2}{16},\qquad BG^2=\frac{7a^2}{16},\qquad CG^2=\frac{7a^2}{16}.
\]
\textbf{3.} Pour tout $M$, $f(M)=4MG^2+f(G)$ et $f(G)=2\cdot\frac{3a^2}{16}+\frac{7a^2}{16}+\frac{7a^2}{16}=\frac{5a^2}{4}$.
Donc $f(M)=3a^2\iff 4MG^2=\frac{7a^2}{4}\iff MG=\dfrac{a\sqrt7}{4}$.
\textbf{$(E_2)$ est le cercle de centre $G$ et de rayon $\dfrac{a\sqrt7}{4}$.}

%% ===============================================================
\subsection{Exercice 28 : Hom\'eth\'eties}

\textbf{1.} $h:\ x'=2x+8,\ y'=2y-4$ : point fixe $x=-8$, $y=4$ ; \textbf{$h$ est l'homoth\'etie de centre $\Omega(-8;4)$ et de rapport $2$}.
\\ $h'$ est l'homoth\'etie de centre $A(-8;4)$ et de rapport $4$ : $\ x'=4x+24,\ y'=4y-12$.
(Tous deux ont le m\^eme centre.)

\textbf{2.} $h\circ h'$ : $(x,y)\mapsto(4x+24,4y-12)\mapsto(2(4x+24)+8,\ 2(4y-12)-4)$, soit
\[
(h\circ h')(x,y)=\bigl(8x+56,\ 8y-28\bigr).
\]

\textbf{3.} Le rapport vaut $2\times4=8$ et le point fixe est $(-8;4)$ (v\'erification : $-8=8(-8)+56$ et $4=8\cdot4-28$).
\textbf{$h\circ h'$ est l'homoth\'etie de centre $(-8;4)$ et de rapport $8$.}

%% ===============================================================
\subsection{Exercice 29 : Transformations}

\textbf{Partie A.} $f:\ x'=3x-4,\ y'=3y+2$.

\textbf{1.} Point invariant : $x=3x-4\iff x=2$ ; $y=3y+2\iff y=-1$. \textbf{$\Omega(2;-1)$} est l'unique point invariant.

\textbf{2.} $\vect{\Omega M'}=3\,\vect{\Omega M}$.

\textbf{3.} $f$ est l'\textbf{homoth\'etie de centre $\Omega(2;-1)$ et de rapport $3$}.

\textbf{Partie B.} $G_m=\mathrm{bar}\{(A;1),(B;m^2),(C;1)\}$.

\textbf{1.} $G_m$ existe si la somme des poids est non nulle : $1+m^2+1=m^2+2\ge2>0$. \textbf{Pour tout $m\in\R$.}

\textbf{2.(a)} Soit $I$ le milieu de $[AC]$ : en vecteurs de position, $G_m=\dfrac{A+C+m^2B}{m^2+2}=\dfrac{2I+m^2B}{m^2+2}$.
C'est le barycentre de $\{(B;m^2),(I;2)\}$ : donc $G_m\in(BI)$.

\textbf{2.(b)} $\vect{BG_m}=\dfrac{\vect{BA}+\vect{BC}}{m^2+2}=\dfrac{2\vect{BI}}{m^2+2}$ (car $\vect{BA}+\vect{BC}=2\vect{BI}$).

\textbf{2.(c)} Pour $m\in\R$, $k=\dfrac{2}{m^2+2}$ d\'ecrit $]0;1]$ :
\begin{itemize}
  \item $m=0$ donne $k=1$, donc $G_0=I$ ;
  \item $|m|\to+\infty$ donne $k\to0$, donc $G_m\to B$ sans jamais l'atteindre.
\end{itemize}
\textbf{L'ensemble des points $G_m$ est le segment $]BI]$} (sans $B$, avec $I$).

\section{Chapitre 5 : Trigonom\'etrie}

%% ===============================================================
\subsection{Exercice 30 : \'Equations trigonom\'etriques}

\textbf{A) QCM.}

\textbf{1.} On a $\cos x-\sqrt3\sin x=2\cos\left(x+\dfrac{\pi}{3}\right)$. L'\'equation $\cos\left(x+\frac{\pi}{3}\right)=0$ donne
$x=\dfrac{\pi}{6}+k\pi$. Dans $]-\pi;\pi]$ : $S=\left\{\dfrac{\pi}{6};-\dfrac{5\pi}{6}\right\}$.
Aucune proposition n'est exacte : \textbf{r\'eponse d)}. (La proposition c) $\{\frac{\pi}{6}\}$ est incompl\`ete.)

\textbf{2.} On retranche un multiple de $2\pi$ : $\dfrac{498\pi}{5}-50\times2\pi=\dfrac{498\pi-500\pi}{5}=-\dfrac{2\pi}{5}$, qui est bien dans $]-\pi;\pi]$.
\textbf{R\'eponse a)} $-\dfrac{2\pi}{5}$.

\textbf{3.} $\cos x=\dfrac{\sqrt6+\sqrt2}{4}=\cos\dfrac{\pi}{12}$, avec $x\in\left[0;\dfrac{\pi}{2}\right]$ : $x=\dfrac{\pi}{12}$.
Alors $\cos2x=\cos\dfrac{\pi}{6}=\dfrac{\sqrt3}{2}$ : \textbf{r\'eponse d)}.

\textbf{B) Vrai ou faux.}
\begin{enumerate}
  \item \textbf{Faux.} Une application surjective n'est pas forc\'ement injective : par exemple $x\mapsto x^2$ de $\R$ vers $\R_+$ est surjective mais pas injective.
  \item \textbf{Vrai.} $\cos x=-\frac12$ sur $]-\pi;\pi]$ donne $x=\pm\frac{2\pi}{3}$, soit $S=\left\{-\frac{2\pi}{3};\frac{2\pi}{3}\right\}$.
\end{enumerate}

%% ===============================================================
\subsection{Exercice 31 : \'Equations et in\'equations}

Sur $[0;2\pi]$ : $\sqrt3\cos2x-\sin2x=2\cos\left(2x+\dfrac{\pi}{6}\right)$. L'in\'equation (I) devient
\[
2\cos\left(2x+\tfrac{\pi}{6}\right)+1<0\iff\cos\left(2x+\tfrac{\pi}{6}\right)<-\frac12.
\]
Posons $\theta=2x+\frac{\pi}{6}$, avec $\theta\in\left[\frac{\pi}{6};\frac{25\pi}{6}\right]$. On a $\cos\theta<-\frac12$ sur
$\left]\frac{2\pi}{3}+2k\pi;\frac{4\pi}{3}+2k\pi\right[$.
\begin{itemize}
  \item $k=0$ : $\theta\in\left]\frac{2\pi}{3};\frac{4\pi}{3}\right[$, soit $2x\in\left]\frac{\pi}{2};\frac{7\pi}{6}\right[$, donc $x\in\left]\frac{\pi}{4};\frac{7\pi}{12}\right[$.
  \item $k=1$ : $\theta\in\left]\frac{8\pi}{3};\frac{10\pi}{3}\right[$, soit $2x\in\left]\frac{5\pi}{2};\frac{19\pi}{6}\right[$, donc $x\in\left]\frac{5\pi}{4};\frac{19\pi}{12}\right[$.
\end{itemize}
\textbf{$S=\left]\frac{\pi}{4};\frac{7\pi}{12}\right[\ \cup\ \left]\frac{5\pi}{4};\frac{19\pi}{12}\right[$} : \textbf{r\'eponse a)}
(\`a l'ouverture pr\`es des intervalles, qui sont ouverts car l'in\'egalit\'e est stricte).

%% ===============================================================
\subsection{Exercice 32 : \'Equations trigonom\'etriques}

\textbf{(E1)} $\cos x-\sin x=\sqrt2$. Comme $\cos x-\sin x=\sqrt2\cos\left(x+\frac{\pi}{4}\right)$ : $\cos\left(x+\frac{\pi}{4}\right)=1$,
donc $x=-\dfrac{\pi}{4}+2k\pi$, $k\in\mathbb Z$.

\textbf{(E2)} $\sqrt3\sin x+3\cos x=\sqrt3$. Comme $\sqrt3\sin x+3\cos x=2\sqrt3\sin\left(x+\frac{\pi}{3}\right)$ :
$\sin\left(x+\frac{\pi}{3}\right)=\frac12$, donc $x+\frac{\pi}{3}=\frac{\pi}{6}+2k\pi$ ou $\frac{5\pi}{6}+2k\pi$.
\[
S_{E_2}=\left\{-\frac{\pi}{6}+2k\pi\ ;\ \frac{\pi}{2}+2k\pi\ ;\ k\in\mathbb Z\right\}.
\]

\textbf{(E3)} $\sqrt3\sin x+\cos x=-1$. Comme $\sqrt3\sin x+\cos x=2\sin\left(x+\frac{\pi}{6}\right)$ :
$\sin\left(x+\frac{\pi}{6}\right)=-\frac12$, donc $x+\frac{\pi}{6}=-\frac{\pi}{6}+2k\pi$ ou $\frac{7\pi}{6}+2k\pi$.
\[
S_{E_3}=\left\{-\frac{\pi}{3}+2k\pi\ ;\ \pi+2k\pi\ ;\ k\in\mathbb Z\right\}.
\]
(V\'erification : pour $x=\pi$, $\sqrt3\cdot0+(-1)=-1$.)

%% ===============================================================
\section{Chapitre 6 : G\'eom\'etrie dans l'espace}

%% ===============================================================
\subsection{Exercice 33 : Cube et orthogonalit\'e}

On utilise le rep\`ere $\left(A;\vect{AB},\vect{AD},\vect{AE}\right)$ avec le cube de c\^ot\'e $a$ :
$A(0;0;0)$, $B(a;0;0)$, $C(a;a;0)$, $D(0;a;0)$, $E(0;0;a)$, $F(a;0;a)$, $G(a;a;a)$, $H(0;a;a)$.
(Figure : face $ABCD$ en bas, face $EFGH$ au-dessus, $E$ au-dessus de $A$.)
Les milieux : $I$ milieu de $[BF]$ : $I\left(a;0;\frac a2\right)$ ; $J$ milieu de $[AE]$ : $J\left(0;0;\frac a2\right)$.

\textbf{1.} Figure : voir le sch\'ema ci-dessus (cube, puis $I$ et $J$ sur les ar\`etes $[BF]$ et $[AE]$).

\textbf{2.} $\vect{ED}=(0;a;-a)$ et $\vect{BG}=(0;a;a)$. Leur produit scalaire vaut $a^2-a^2=0$.
\textbf{Donc $(ED)\perp(BG)$.}

\textbf{3.} $\vect{IJ}=\left(-a;0;0\right)$. Le plan $(ADE)$ est le plan $x=0$, de vecteur normal $(1;0;0)$ :
$\vect{IJ}$ est normal au plan. \textbf{Donc $(IJ)\perp(ADE)$.}

\textbf{4.} Plan $(EAC)$ : il contient $A$, $E$ et $C$, donc un vecteur normal est $\vect n_1=(1;-1;0)$
(orthogonal \`a $\vect{AE}=(0;0;a)$ et $\vect{AC}=(a;a;0)$).
Plan $(HDB)$ : $\vect{DB}=(a;-a;0)$ et $\vect{DH}=(0;0;a)$, donc un vecteur normal est $\vect n_2=(1;1;0)$.
$\vect n_1\cdot\vect n_2=1-1=0$. \textbf{Donc les plans $(EAC)$ et $(HDB)$ sont perpendiculaires.}

\textbf{5.} Plan $(DBF)$ : $\vect{DB}=(a;-a;0)$ et $\vect{DF}=(a;-a;a)$. Le vecteur $(1;1;0)$ est normal ($\vect{DB}\cdot(1;1;0)=0$ et
$\vect{DF}\cdot(1;1;0)=0$) ; le plan a pour \'equation $x+y=a$ (v\'erifi\'ee en $D$, $B$, $F$).
\[
d(A,(DBF))=\frac{|0+0-a|}{\sqrt{1^2+1^2}}=\frac{a}{\sqrt2}=\frac{a\sqrt2}{2}.
\]

%% ===============================================================
\section{Chapitre 7 : QCM et questions diverses}

%% ===============================================================
\subsection{Exercice 34 : QCM}

\textbf{A.1.} $f(x)=\dfrac{5x^3-3x-2}{x^4-x^3}$ (lecture de l'\'enonc\'e). Au voisinage de $+\infty$, $f(x)\sim\dfrac{5x^3}{x^4}=\dfrac5x\to0$.
La limite vaut $0$ : \textbf{r\'eponse d)} (aucune des trois autres).
\begin{rem} Si le num\'erateur est diff\'erent, la conclusion peut changer ; v\'erifier l'\'enonc\'e.\end{rem}

\textbf{A.2.} $g(x)=\dfrac{-2x^2-x+1}{x^2+1}$. Le d\'enominateur ne s'annule jamais : \textbf{pas d'asymptote verticale}.
En $\pm\infty$, $g(x)\to-2$ : la droite $y=-2$ est asymptote horizontale. \textbf{R\'eponse b).}

\textbf{B.1.} La tangente \`a $\mathcal C_f$ au point d'abscisse $a$ a pour \'equation $\boxed{y=f'(a)(x-a)+f(a)}$.

\textbf{B.2.} $\displaystyle\lim_{x\to0}\frac{\sin x}{x}=\mathbf{1}$.

\textbf{B.3.} Rotation de centre $O$ et d'angle $\frac{2\pi}{3}$ : $\cos\frac{2\pi}{3}=-\frac12$ et $\sin\frac{2\pi}{3}=\frac{\sqrt3}{2}$, donc
\[
\begin{cases}
x'=-\dfrac12x-\dfrac{\sqrt3}{2}y\\[6pt]
y'=\dfrac{\sqrt3}{2}x-\dfrac12y.
\end{cases}
\]

%% ===============================================================
\subsection{Exercice 35 : Vrai/Faux et compl\'ements}

\textbf{Partie A.}
\begin{enumerate}
  \item \textbf{Vrai.} Une fonction bijective est injective (et surjective).
  \item \textbf{Faux.} $\sin x\in[-1;1]$, donc $\sin x=\frac32$ n'a aucune solution.
\end{enumerate}

\textbf{Partie B.}
\begin{enumerate}
  \item Une fonction $f$ est une \textbf{application} si l'ensemble de d\'epart et l'ensemble de d\'efinition coïncident.
  \item Deux r\'eels de somme $S$ et de produit $P$ sont les solutions de $\boxed{x^2-Sx+P=0}$.
  \item $f:E\to F$ est \textbf{injective} si chaque \'el\'ement de $F$ admet au plus un ant\'ec\'edent.
  \item $x+y=\dfrac{2\pi}{3}-\dfrac{\pi}{4}=\dfrac{5\pi}{12}$, d\'ej\`a dans $]-\pi;\pi]$ : \textbf{mesure principale $\dfrac{5\pi}{12}$}.
  \item $h(M)=M'$ \'equivaut \`a $\boxed{\vect{\Omega M'}=-2\,\vect{\Omega M}}$.
  \item Si $G=\mathrm{bar}\{(A;\alpha);(B;\beta)\}$ avec $\alpha+\beta=1$, alors $\vect{BG}=\alpha\,\vect{BA}$.
  Or $\vect{BG}=-\dfrac23\vect{BA}$ donne $\alpha=-\dfrac23$ et $\beta=1-\alpha=\dfrac53$.
\end{enumerate}

%% ===============================================================
\subsection{Exercice 36 : QCM}

\textbf{I. Vrai ou faux.}
\begin{enumerate}
  \item \textbf{Vrai.} $f(x)=\sin\left(2x-\frac{\pi}{3}\right)$ a pour p\'eriode $\frac{2\pi}{2}=\pi$.
  \item \textbf{Faux.} $\cos x=\frac14$ et $\cos y=\frac12$ avec $x,y\in\left[0;\frac{\pi}{2}\right]$ : $y=\frac{\pi}{3}$, $\sin x=\frac{\sqrt{15}}{4}$.
  $\cos2y=-\frac12$, $\sin2y=\frac{\sqrt3}{2}$ et
  \[
  \cos(x-2y)=\cos x\cos2y+\sin x\sin2y=-\frac18+\frac{\sqrt{45}}{8}=\frac{3\sqrt5-1}{8}.
  \]
  L'\'enonc\'e donne $\frac{3\sqrt5+1}{8}$ : \textbf{l'affirmation est fausse}.
  \item \textbf{Faux.} $\vect{AB}=3\vect{CA}$ donne $\vect{AB}=-3\vect{AC}$ : $B$ est l'image de $C$ par l'homoth\'etie de centre $A$ et de rapport $-3$ (et non $3$).
  \item \textbf{Vrai.} $MA^2-MB^2=0\iff MA=MB$ : c'est la m\'ediatrice de $[AB]$.
\end{enumerate}

\textbf{II. Choix multiple} (\'enonc\'e non reproduit dans le texte).

%% ===============================================================
\subsection{Exercice 37 : QCM}

\textbf{I.1.} $3x^2-2|x|-1=0$. En posant $t=|x|\ge0$ : $3t^2-2t-1=0$, soit $(3t+1)(t-1)=0$, donc $t=1$ (car $t\ge0$).
Alors $x=\pm1$ : \textbf{deux solutions distinctes, r\'eponse b)}.

\textbf{I.2.} $\sqrt{x^2-2x}>1-x$ : domaine $x^2-2x\ge0$, soit $x\le0$ ou $x\ge2$.
\begin{itemize}
  \item Si $1-x<0$ (donc $x>1$), l'in\'egalit\'e est vraie d\`es que $x\ge2$.
  \item Si $1-x\ge0$ et $x\le0$ : on \'el\`eve au carr\'e, $x^2-2x>(1-x)^2=1-2x+x^2$, soit $0>1$ : impossible.
\end{itemize}
Donc $S=[2;+\infty[$ : \textbf{r\'eponse a)}.

\textbf{I.3.} L'\'enonc\'e \'ecrit $3\vect{BC}-4\vect{AC}=2\vect{BC}$, soit $\vect{BC}=4\vect{AC}$, ce qui s'\'ecrit $4A-B-3C=\vec 0$.
Cette relation n'est pas de la forme $\sum\alpha_i\vect{GX_i}=\vec 0$ avec les poids des options (les poids de l'\'enonc\'e ne somment pas \`a $0$) :
\begin{rem} L'\'enonc\'e est incoh\'erent (il manque probablement un point ou une lettre). Aucune des quatre r\'eponses ne se d\'eduit
directement de la relation \'ecrite ; merci de v\'erifier la formule originale.\end{rem}

%% ===============================================================
\subsection{Exercice 38 : Ajustement lin\'eaire}

\textbf{Consigne 1.} La production journali\`ere est $h(v)=-v+3300-\dfrac{8100}{v}$ pour $v\in[50;120]$.
\[
h'(v)=-1+\frac{8100}{v^2}=\frac{8100-v^2}{v^2}.
\]
$h'(v)=0\iff v=90$ (on prend $v>0$) ; $h'>0$ avant, $h'<0$ apr\`es : maximum en $v=90$.
$h(90)=-90+3300-90=3120$. \textbf{La production maximale est atteinte pour $v=90$, avec $3120$ exemplaires par jour.}

\medskip\noindent\textbf{Consigne 2.} Ajustement par moindres carr\'es de $y$ en fonction de $x$ ($x$ en milliers de FCFA) :
$\bar x=130$, $\bar y=504$, $\sum(x_i-\bar x)^2=16\,800$, $\sum(x_i-\bar x)(y_i-\bar y)=-100\,060$.
\[
a=\frac{-100\,060}{16\,800}\approx-5{,}956,\qquad b=\bar y-a\bar x\approx1\,278{,}27,\qquad \hat y=-5{,}956\,x+1\,278{,}27.
\]
\textbf{B\'en\'efice} : chaque exemplaire est vendu $1000x$ F et co\^ute $25\,000$ F ; les frais fixes sont $28\,000\,000$ F. Donc
\[
B(x)=\hat y\,(1000x-25\,000)-28\,000\,000.
\]
En d\'eveloppant : $B(x)=-5\,956\,x^2+1\,427\,170\,x-59\,956\,750$.
Le maximum est atteint pour $x=\dfrac{1\,427\,170}{2\times5\,956}\approx119{,}8$.
\textbf{Le prix de vente optimal est d'environ $119\,800$ F, pour un b\'en\'efice maximal d'environ $25{,}5$ millions de FCFA.}
\begin{hyp}le mod\`ele lin\'eaire $\hat y(x)$ est utilis\'e pour pr\'evoir le nombre d'exemplaires vendus.\end{hyp}

%% ===============================================================
\subsection{Exercice 39 : Probabilit\'es}

Une classe de $12$ filles et $9$ gar\c cons, soit $21$ \'el\`eves. Un comit\'e est un arrangement de $3$ \'el\`eves
(major, adjoint, secr\'etaire).

\textbf{1.} Nombre de comit\'es : $21\times20\times19=\mathbf{7980}$ : \textbf{r\'eponse b)}.

\textbf{2.} Comit\'es de m\^eme sexe : tous filles $12\times11\times10=1320$ ; tous gar\c cons $9\times8\times7=504$.
Total : $1320+504=\mathbf{1824}$ : \textbf{r\'eponse c)}.

\textbf{3.} $P(\text{que des filles})=\dfrac{1320}{7980}=\dfrac{22}{133}\approx0{,}165$ : \textbf{r\'eponse b)}.

%% ===============================================================
\subsection{Exercice 40 : Lancers de pi\`ece}

On lance deux fois une pi\`ece : $\Omega=\{(F;F);(F;P);(P;F);(P;P)\}$. Soit $A$ : ``obtenir au moins un pile'',
donc $A=\{(F;P);(P;F);(P;P)\}$.

\textbf{Appariement :}
\begin{itemize}
  \item ``Univers'' $\longleftrightarrow$ $\Omega=\{(F;F);(F;P);(P;F);(P;P)\}$ ;
  \item ``Obtenir uniquement les c\^ot\'es face'' $\longleftrightarrow$ l'\'ev\'enement $\{(F;F)\}$, qui est un \textbf{\'ev\'enement \'el\'ementaire} ;
  \item ``\'Ev\'enement incompatible avec $A$'' $\longleftrightarrow$ $\{(F;F)\}$, car $A\cap\{(F;F)\}=\varnothing$ ;
  \item ``\'Ev\'enement \'el\'ementaire'' $\longleftrightarrow$ $(P;P)$ ;
  \item ``\'Eventualit\'e'' $\longleftrightarrow$ $(F;P)$.
\end{itemize}
\begin{rem} L'\'enonc\'e propose une liste d'appariements dont l'un se recoupe : ``obtenir uniquement face'' est \`a la fois
\'ev\'enement \'el\'ementaire et incompatible avec $A$. La r\'eponse ci-dessus est la plus coh\'erente.\end{rem}

%% ===============================================================
\subsection{Exercice 41 : Cercle et repr\'esentation param\'etrique}

$x^2+y^2-12x+6y-4=0\iff(x-6)^2+(y+3)^2=4+36+9=49$.
\textbf{Centre $\Omega(6;-3)$, rayon $R=7$.}
\textbf{Repr\'esentation param\'etrique :} $\begin{cases}x=6+7\cos t\\ y=-3+7\sin t\end{cases}$, $t\in[0;2\pi[$.

%% ===============================================================
\subsection{Exercice 42 : Droites et angles}

Pour deux droites $(\Delta)$ et $(\Delta')$ s\'ecantes en $A$, la compos\'ee de deux sym\'etries d'axes concourants est une rotation de centre $A$ :
\[
S_{(\Delta')}\circ S_{(\Delta)}=\text{rotation de centre }A\text{ et d'angle }2\,\mathrm{Mes}(\vect u;\vect u')
=2\times\left(-\frac{\pi}{6}\right)=-\frac{\pi}{3}.
\]
\textbf{R\'eponse : la rotation de centre $A$ et d'angle $-\dfrac{\pi}{3}$.}

%% ===============================================================
\subsection{Exercice 43 : Compos\'ee de sym\'etries}

\textbf{La compos\'ee de deux sym\'etries d'axes parall\`eles est une \textbf{translation}} ; par contre la compos\'ee de deux
sym\'etries d'axes perpendiculaires est une \textbf{sym\'etrie centrale} (rotation d'angle $\pi$) de centre le point d'intersection des axes.

\section*{Conclusion}
Ces corrections couvrent les 43 exercices du recueil. Les points marqu\'es ``remarque'' ou ``hypoth\`ese'' signalent des
\'enonc\'es incomplets, incoh\'erents ou ambigus : il est conseill\'e de les v\'erifier aupr\`es du professeur avant de les utiliser.

\end{document}
